\SourceLeaf{412}%
\par
So hat man folgende Relationen
\[
\left.\begin{aligned}
tu.\sin p&=2a\\ tv.\sin q&=2x\\ tw.\sin r&=2e\\
vw.\sin(r-q)&=2d\\ uw.\sin(r-p)&=2z\\ uv.\sin(q-p)&=2y
\end{aligned}\right\}\quad
\begin{gathered}\text{daraus vermittelst des Lemma der}\\
\textit{Th. M. C. C.}\\ ad-xz+ey=0\qquad\text{(I)}\end{gathered}
\]
\[
\left.\begin{aligned}b+d+x&=\omega\\a+d+y&=\omega\\a+c+z&=\omega\end{aligned}\right\}\quad
\begin{gathered}\text{die Werthe von }x,y,z\text{ hieraus in (I)}\\\text{substituirt geben}\end{gathered}
\]
\[ad-(\omega-b-d)(\omega-a-c)+e(\omega-a-d)=0\]
oder entwickelt
\[\omega\omega-(a+b+c+d+e)\omega+(ab+bc+cd+de+ea)=0\]
\begin{flushright}\textsc{Schumacher}.\}\end{flushright}
\par\bigskip\Needspace{7\baselineskip}
\pdfbookmark[1]{Handbuch der Schiffahrtskunde. 1850.}{schiffahrtskunde-1850}
\begin{center}Handbuch der Schiffahrtskunde von C. \textsc{Rümker}. 1850. Seite 76.\par\medskip
\textit{Auflösung einer geometrischen Aufgabe.}\end{center}
\hspace*{\parindent}An drei Punkten \((1),(2),(3)\), welche in einer geraden Linie \(\mathrm{(I)}\) und in bekannten Abständen von einander, \(A\) von \((1)\) nach \((2)\), \(B\) von \((2)\) nach \((3)\) liegen, sind die Winkel \(\vartheta,\vartheta',\vartheta''\) zwischen zwei andern Punkten \((4),(5)\), deren gegenseitiger Abstand \(=2c\) ebenfalls bekannt ist, gemessen; man verlangt die Lage der drei erstern Punkte gegen die beiden letztern. Um nichts unbestimmt zu lassen, setze ich voraus, dass die drei Winkel alle von \((4)\) nach \((5)\) in einerlei Sinn wachsend gemessen sind, dass auf der Linie \(\mathrm{(I)}\) die Abstände in einem bestimmten Sinne positiv gezählt werden (so dass, wenn man aus irgend welchem Grunde nicht den zwischen den beiden andern liegenden Punkt mit \((2)\) bezeichnete, \(A\) und \(B\) ungleiche Zeichen erhalten würden) und \(c\) positiv genommen werden soll.
